% Section 6 — Learning, contribution, practical implications
\section{Learning, Contribution and Practical Implications}
\label{sec:contrib}

\subsection{Empirical contribution}
\label{sec:c-emp}
The study measures the forecast-decision relationship in both directions
within one controlled design: correspondence on intermittent M5 demand and
divergence on dense Store demand, with the Store podium (Moving Average,
SES, Seasonal Naive ahead of LSTM) as the sharpest case. It records where
sophistication pays (LSTM accuracy on both datasets; LSTM cost on sparse
demand) and where it does not (TES and DES; ARIMA on dense demand).

\subsection{Methodological contribution}
\label{sec:c-meth}
Three reusable instruments. First, a shared-policy evaluation harness that
holds the inventory treatment constant across models: one module, every
model, no per-model tuning. Second, a leakage-safe rolling protocol (history
strictly before origin, validation-only selection, honest missing-value
handling in scale-free metrics with pooled archetype ratios). Third, a
stability discipline in which headline rankings are published together with
their 27-policy record, separating stable conclusions from conditional ones.

\subsection{Analytical contribution}
\label{sec:c-anal}
Ninety-one Holm-corrected paired comparisons with effect sizes and bootstrap
intervals, reported under the rule that $p$-values do not travel without
$d_z$. The MA-SES reference null shows the procedure returning equivalence
as well as difference.

\subsection{Inventory-management implications}
\label{sec:c-inv}
For intermittent catalogs, learned or specialist forecasting (LSTM, TSB/SBA)
is supported, and the forecast advantage reaches the shelf. For dense,
fast-moving catalogs, simple averaging methods should be piloted against the
neural option under the applicable cost ratio before engineering investment;
the lowest-cost input may already be in production. In both settings,
leaderboards that stop at accuracy are insufficient evidence for procurement.

\subsection{Model-selection implications}
\label{sec:c-select}
Selection is conditional on three answers: the demand structure (sparse or
dense), the cost asymmetry ($P/H$), and the lead time. The report's tables
cover the remainder: sparse demand with standard policies favors LSTM; dense
demand under the default policy favors Moving Average with SES as the
alternative; dense demand with short leads and high penalties requires a
fresh comparison, since leadership is fluid there.

\subsection{Reproducibility contribution}
\label{sec:c-repro}
The frozen package (number sheet, audits, traceability, validation log with
0 FAIL, pre-hardening backup, seeds, manifests) maps every reported number
to a source file in one lookup. Provenance details are in \cref{app:repro}.

\subsection{Research-hardening lessons}
\label{sec:c-harden}
The costliest errors found were methodological rather than algorithmic: an
unassigned leap day, overlapping validation windows, a subset result carried
at full weight, and a mean-of-ratios construction unstable on near-zero
series. Each was identified by audits that assumed nothing (row counts, date
containment, denominator discipline, population labels). The general lesson
is that the experiment requires the same scrutiny as the model, since
conclusions are drawn from the former.
